\section{The complete two-counter highway basin}\label{sec:basin}
A standard highway is permanent
agreement of the pose and turn trace with the usual translating period-$104$
word, up to phase and quarter-turn rotation. Historical black cells need not
become a translated finite motif. To fix the convention completely, our standard
word is the concatenation of these two displayed strings:
\begin{center}\scriptsize\ttfamily
RRLLLLRLRRRRLRRRRLLLLRLRRRRLRLLRRLLLLRRLLRRRRLLRRLRL\\
LRLLRLRLLRRRRLLRLLRRRRLLRRRRLLRLRRRRLRLLLLRRRRLRRLRR
\end{center}
It has primitive length 104; its cyclic phases and quarter-turns are allowed.

For the remainder of the paper we translate the source of
Section~\ref{sec:five} so its leftmost column is $x=0$. This changes its
initial head from $(0,3)$ to $(2l+5,3)$ and leaves all durations unchanged.
\subsection{Encoding}
A column of height five is the integer $\sum_{y=0}^4b_y2^y$. Column words are
placed from $x=0$ to the right in rows $0,\ldots,4$; every omitted lattice cell
is white. Set
\[
 A=(6,5),\quad B=(10,9),\quad C=(6,10,12,17),\quad V=(18,29),
 \qquad \R(j,k)=A^jCV^k.
\]
For a column word define
\[
 \J(c_1,\ldots,c_n)=(31-c_n,\ldots,31-c_1).
\]
Thus $\J$ reflects and complements \emph{inside the indicated finite block};
it does not complement the exterior. For $l,r\in\N=\{0,1,\ldots\}$, let
$X(l,r)$ have column word and initial pose
\begin{equation}
 \J((5)CV^l)\,CV^r,\qquad ((2l+5,3),\E).                 \label{basin:eq:source}
\end{equation}
Its width is $n=2(l+r)+9$, its initial black mass is $4l+6r+23$, and all these
cells are retained throughout our argument.

\begin{theorem}[Complete counter basin]\label{basin:thm:main}
For every $l,r\in\N$, the source $X(l,r)$ eventually enters the standard
highway. Its earliest permanent entrance time is
\[
 \tau(l,r)=E(l,r)+\rho(l,r),\qquad
 E(l,r)=\begin{cases}8r^2+60r+16,&r\le l,\\8l^2+68l+48,&l<r,\end{cases}
\]
where $E$ is the first rectangle-exit time and $\rho$ is completely specified
by Tables~\ref{basin:tab:recovery} and~\ref{basin:tab:small}. In the balanced case,
\begin{equation}
 \tau(k,k)=8k^2+100k+1047\quad(k\ge2),                    \label{basin:eq:balanced}
\end{equation}
while $\tau(0,0)=11576$ and $\tau(1,1)=1024$.
\end{theorem}

\begin{table}[t]
\centering
\begin{tabular}{lll}
\toprule
Region & Condition & Recovery delay $\rho$\\\midrule
$r\le l$, $d=l-r$ & $d=0,r=0$ & $11560$\\
 & $d=0,r=1$ & $940$\\
 & $d=0,r\ge2$ & $40r+1031$\\
 & $d=1$ & $32r+883$\\
 & $d\ge2$ & $32r+799$\\\midrule
$l<r$, $l\ge9$, $d=r-l$ & $d=1$ & $72l+1427$\\
 & $d=2$ & $128l+2611$\\
 & $d=3$ & $144l+3011$\\
 & $d=4$ & $176l+3658$\\
 & $d=5$ & $176l+3678$\\
 & $d\ge6$ & $176l+3570$\\\bottomrule
\end{tabular}
\caption{All unbounded parameter cones. The remaining fixed-$l$ strips are
covered by Table~\ref{basin:tab:small}, not omitted.}\label{basin:tab:recovery}
\end{table}

\begin{table}[t]
\centering\small
\begin{tabular}{rl}
\toprule
$l$ & Delays for $d=r-l=1,2,\ldots$; the final entry repeats forever\\\midrule
0 & $17106,2789,1479,2661,5711,2665,2557$\\
1 & $1516,1423,1591,6410,6430,6322$\\
2 & $1547,1925,1945,1837$\\
3 & $1635,2021,2041,1933$\\
4 & $1715,2883,3787,3915,3413,3433,3325$\\
5 & $1787,3071,3447,4199,4495,4859,6975,7638,7658,7550$\\
6 & $1859,4607,4451,4471,4363$\\
7 & $1931,3543,3931,4491,4771,4470,4490,4382$\\
8 & $2003,3635,4163,6057,6077,5969$\\\bottomrule
\end{tabular}
\caption{The finite base table and its protected infinite extensions. Each row
is a formula for all $d\ge1$, not an observation up to its printed length.}
\label{basin:tab:small}
\end{table}

The theorem is computer-assisted in a precise finite sense. The infinite
domains are checked by integer-affine identities with proved coverage, and
only fifty-nine isolated states require literal finite execution. We supply
full words, source formulas and exact verification programs. None of the
parameter cutoffs in the finite controls is a premise of the theorem.

The round induction and two terminal kernels in
Section~\ref{sec:five}, after the indicated translation, give the complete exit states
\begin{equation}
\begin{array}{ll}
r\le l:&\J((5)\R(r,l-r))B^{r+1}(10,18),\quad (n,1),\E,\\
l<r:&\J((9)B^{l+1}(10,18))\R(l+1,r-l-1),\quad (-1,1),\West.
\end{array}                                                    \label{basin:eq:exit}
\end{equation}
\subsection{Recovery certificates for every parameter}
\subsubsection{The nine affine-word programs}
A program is a finite concatenation
\begin{equation}
 u_0v_1^{a_1k+b_1}u_1\cdots v_s^{a_sk+b_s}u_s,             \label{basin:eq:program}
\end{equation}
where the $u_i,v_i$ are explicit LR words. Each variable block $v_i$ restores
the heading and has displacement $(2,0)$ or $(-2,0)$ in its actual pose. The
counts are positive on the declared entire ray $k\ge k_0$.
Table~\ref{basin:tab:schemas} gives the domains and sizes. The literal fixed words
and integer exponents are the immutable data in \word{witness.json}.

\begin{table}[t]
\centering\small
\begin{tabular}{llllrr}
\toprule
Program & $(l,r)$ & Domain & Length & Powers & Read tests\\\midrule
E0 & $(k,k)$ & $k\ge2$ & $40k+1031$ & 6 & 1007\\
E1 & $(k+1,k)$ & $k\ge0$ & $32k+883$ & 4 & 839\\
E2 & $(k+2,k)$ & $k\ge0$ & $32k+799$ & 4 & 751\\
W1 & $(k,k+1)$ & $k\ge9$ & $72k+1427$ & 9 & 1431\\
W2 & $(k,k+2)$ & $k\ge9$ & $128k+2611$ & 17 & 2795\\
W3 & $(k,k+3)$ & $k\ge9$ & $144k+3011$ & 19 & 3203\\
W4 & $(k,k+4)$ & $k\ge9$ & $176k+3658$ & 22 & 4050\\
W5 & $(k,k+5)$ & $k\ge9$ & $176k+3678$ & 22 & 4074\\
W6 & $(k,k+6)$ & $k\ge9$ & $176k+3570$ & 22 & 3970\\\bottomrule
\end{tabular}
\caption{A test for a variable phase covers all its repetitions, not one
sample. The sum is 22,120 read obligations on unbounded domains.}\label{basin:tab:schemas}
\end{table}

For illustration, E0 has fixed-block lengths
$114,108,3,9,3,169,577$, separated by powers of
\[
\begin{split}
 &\word{RLLR}^{2k+3},\quad\word{LRRL}^{2k+3},\quad\word{RRLL}^{2k+3},\\
 &\word{LRRL}^{k+1},\quad\word{RRLL}^{k+1},\quad\word{LRRL}^{2k+1}.
\end{split}
\]
The length is $40k+1031$. This word pattern is not accepted merely because
finite runs exhibit it: every read is proved correct by the following algebra.

\subsubsection{Complete affine fields}
Throughout a program the row of a phase is independent of $k$ and its horizontal
position is $ak+b$, with $a$ even. Variable repetitions add $\pm2j$. Fix a row
$y$ and horizontal parity $p\in\{0,1\}$, and write $x=2z+p$. A black arithmetic interval
is represented exactly by two jumps:
\begin{equation}
 [L(k)\le z\le U(k)]
   =[z\ge L(k)]\mathbin\oplus[z\ge U(k)+1].               \label{basin:eq:jumps}
\end{equation}
Here $L,U$ have integer affine coefficients. Equal jumps cancel by XOR.
Zero-length intervals cancel as well. This representation includes all
initial cells and all complete earlier writes; it never extracts an assumed
active subset.

Suppose the phase reads $q_i+jd$ in a variable block, with $d=(\pm2,0)$.
The pre-block field is an XOR of tests $j\ge f(k)$ for integer-affine $f$.
For a previous phase $q_t$, the unique copy that could have written this cell
has index $j-s$, where $q_t-q_i=sd$. It is earlier exactly when
$s>0$, or $s=0$ and $t<i$, and it exists exactly when $j\ge s$.
Include those constant thresholds and the required L bit at $j=0$. The read
condition becomes
\begin{equation}
 \bigoplus_{f\in\mathcal E}[j\ge f(k)]=0,
 \qquad k\ge k_0,\quad 0\le j<ak+b.                     \label{basin:eq:read}
\end{equation}
Fixed reads are the special case with one epoch. Once a block is verified,
all of its phase intervals are added to the complete field. Induction over
ordinary time proves that satisfying these constraints executes the literal
program, rather than just an abstract trace.

\begin{lemma}[Exact affine epoch test]\label{basin:lem:affine}
For finite $\mathcal E$ and a positive affine epoch count on $k\ge k_0$,
condition~\eqref{basin:eq:read} is decidable by finite integer sign and intersection
checks, including the final infinite parameter interval.
\end{lemma}
\begin{proof}
First check parity at $j=0$. It changes only at the roots $f(k)=0$; isolating
$\lfloor\text{root}\rfloor$ and its successor separates all integer sign regions.
For every distinct affine jump $f$, solve
$1\le f(k)\le ak+b-1$. This is an integer interval. If it is infinite, the
jump cannot coincide with other distinct affine jumps throughout it: finitely
many other lines give only finitely many intersections. Thus the condition
fails. If the interval is finite, enumerate every integer intersection with
other jumps. Their distinct parameter values must cover the whole interval,
and every coincident jump there must have even multiplicity. These conditions
are exactly absence of odd interior jumps. Along with zero initial parity they
prove the claimed XOR identity.
\end{proof}

A second implementation instead partitions at every pairwise affine-order
root, including comparisons with zero and the epoch count. In each resulting
integer interval all order relations are fixed, so one representative checks
initial parity and every interior multiplicity. Both algorithms use exact
integers; neither uses floating-point feasibility, a parameter cutoff, or an
SMT solver. They agree on the complete fields and lengths of all nine programs.
The common geometric compilation remains an explicitly stated part of the
verification boundary.

\subsection{Permanent standard tails and unused data}
\subsubsection{An all-future phase test}
Let $B(k)$ be a program's complete endpoint field. For a proposed standard word,
write $q_i(k)$ for its 104 phase reads and $d$ for its diagonal displacement.
At phase $i$, epoch $n$, its actual read bit is
\begin{equation}
 \mathbf{1}_{B(k)}(q_i(k)+nd)
 \oplus\#\{t:q_t-q_i=sd,\ s>0\text{ or }(s=0,t<i),\ s\le n\}\pmod2.
                                                               \label{basin:eq:tail}
\end{equation}
This is obtained directly by counting all preceding planned flips. Because
$B(k)$ uses finitely many rows independently of $k$ and $d_y=\pm2$, a phase
ray meets each such row at only one possible nonnegative epoch. Check those
epochs and their successors, all prior-write thresholds, zero and one epoch
beyond the last change. Between them the bit is constant. At each chosen epoch
the remaining condition is an affine test in $k$ handled by
Lemma~\ref{basin:lem:affine}. This proves all epochs, not a fixed number of repetitions.

There are 4,415 such all-parameter phase/epoch checks. The certified word is a
cyclic phase of the standard primitive word. E0--E2 and W1--W3 finish at
$(n+3,7),\North$; W4--W6 finish at $(14,14),\West$. Their tail drift is
$(-2,2)$. The standard-word identity is also independently checked against a
complete blank-orbit state and the same infinite criterion.

\begin{lemma}[Earliest entrance]\label{basin:lem:first}
Each listed completion time is the earliest permanent standard time.
\end{lemma}
\begin{proof}
The last actual pre-completion letter differs from the letter immediately
preceding the certified standard phase. An earlier permanent standard suffix
would have the same unique 104-periodic backwards extension, since the standard
word is primitive and both futures agree. This contradicts that letter. The
argument also excludes an earlier permanent suffix during the counter phase.
\end{proof}

\subsubsection{Protection of arbitrary extra counter lengths}
E2 proves the base $l=r+2$. Increasing $l$ prepends more copies of $\J(V)$.
After aligning the common part, the difference lies in $x<0$, rows $0,\ldots,4$.
All finite-program reads in those rows have $x\ge0$. For each variable scan,
it is enough to check both affine endpoints; all future phase-ray encounters
with those rows are checked as well. Consequently the whole actual future
ignores that difference, by induction up to a hypothetical first differing
read. The added cells stay in the field.

Similarly W6 extends from $r=l+6$ to every larger $r$: the extra $V$ blocks
are beyond $x=4l+21$ in the five data rows, and no finite-program read or future
phase-ray encounter reaches them. The final finite base in each row of
Table~\ref{basin:tab:small} has the corresponding guard. Some small cases eventually
drift right, but their rays leave the data rows before reaching the extensions;
a drift sign alone is not used as a protection proof.

The two balanced exceptions and all 57 left-exhausted finite bases are replayed
from~\eqref{basin:eq:exit}. Their complete fields satisfy~\eqref{basin:eq:tail} and
Lemma~\ref{basin:lem:first}; the designated nine bases also satisfy the extension
guards. These domains exhaust $\N^2$. Together with the nine programs, this
proves Theorem~\ref{basin:thm:main}.

\begin{corollary}[Exact trace clipping]\label{basin:cor:clip}
For $l\ge r$, replace $l$ by $\min(l,r+2)$. For $l<r$, replace $r$ by
$\min(r,l+D(l))$, where $D(l)$ is the length of row $l$ in
Table~\ref{basin:tab:small} for $l<9$, and $D(l)=6$ otherwise. After aligning their
initial heads, the original and clipped inputs have the same entire forward
head and turn trace. Their complete fields differ by the static untouched
extra columns. In particular $D(l)\le10$ for every $l$.
\end{corollary}
\begin{proof}
During the initial counter exchanges the six kernels cannot reach the omitted
blocks: both counters make only the shorter number of successful exchanges,
with one additional right success when the left counter exhausts first.
The complete post-exit states therefore differ only in those blocks. The
finite and infinite guards just proved continue the equality forever.
\end{proof}

The same first-disagreement argument allows arbitrary finite payloads in the
protected omitted half-strips, rather than only the indicated counter pattern.
This is a proved guard for these specific common sources, not a procedure for
discarding arbitrary cells by assumed future liveness.

\subsection{Exact extrema and nonlinear entrance times}
\begin{theorem}[Worst entrance at fixed smaller counter]\label{basin:thm:max}
The quantity $M(m)=\max_{\min(l,r)=m}\tau(l,r)$ is finite.
The maximum allows unbounded width and aspect ratio.
For $m=0,\ldots,8$, respectively, its values are
\[
 17154,6554,2161,2365,4363,8246,5351,5687,7181,
\]
and for every $m\ge9$,
\begin{equation}
 M(m)=8m^2+244m+3726.                                     \label{basin:eq:max}
\end{equation}
Consequently, for all $l,r\ge0$ and $m=\min(l,r)$,
\begin{equation}
 \tau(l,r)\le8m^2+244m+17154.                             \label{basin:eq:uniform}
\end{equation}
The coefficients $8,244,17154$ are lexicographically optimal among bounds of
this polynomial form.
\end{theorem}
\begin{proof}
For $m\ge9$, the choice $(l,r)=(m,m+5)$ gives~\eqref{basin:eq:max}. Subtracting the
other affine delays in Table~\ref{basin:tab:recovery} proves that none is larger.
For $m<9$, maximize the three right-exit cases and the finite row of left-exit
values; the last row value already covers all larger aspect ratios. This gives
the displayed table and~\eqref{basin:eq:uniform}. The ray $(m,m+5)$ forces first the
quadratic coefficient to be at least eight, then, with eight fixed, the linear
coefficient to be at least 244. With both fixed, $(l,r)=(0,1)$ forces the constant
to be at least 17154.
\end{proof}

The lower bound $E(l,r)$ and~\eqref{basin:eq:uniform} show, uniformly in the larger
counter,
\[
 \tau(l,r)=8\min(l,r)^2+O(\min(l,r)+1).
\]
This uniformity is a dynamical conclusion: adding an arbitrarily long memory
tail does not necessarily prolong the execution, because a verified finite
margin screens it from every future read.

\begin{corollary}[Nonsemilinear entrance graph]\label{basin:cor:nonsemi}
The complete source family~\eqref{basin:eq:source} is Presburger-definable with
finite fibers. Every member lies in the standard-highway basin, but the graph
$\{(l,r,\tau(l,r)):l,r\in\N\}$ is not semilinear.
\end{corollary}
\begin{proof}
The column boundaries are affine and the two-column words use congruences
modulo two, so the source family is Presburger-definable. If the entrance graph
were semilinear, intersecting with $l=r=k\ge2$ and projecting would give a
semilinear graph for~\eqref{basin:eq:balanced}. An infinite linear component of such
a function graph must contain an arithmetic ray with positive increment $a$
in $k$. Its second difference in the output is zero, whereas the quadratic has
second difference $16a^2>0$. Hence no infinite linear component is possible,
and finitely many finite components cannot cover the graph.
\end{proof}

Every nonnegative parameter pair belongs to the highway basin, and its entrance
time is evaluable by the displayed integer formulas in polynomial bit time.
The correspondence
between Presburger-definable sets and semilinear sets is the classical
theorem of Ginsburg and Spanier~\cite[Theorem~1.3]{gs}.

